\section{Artin環上の標準化と周期的変形複体}\label{sec:deformation-complex}
以下、変形とは局所Artin $k$-代数$B$上の変形であり、特殊ファイバーとの同定を保持する。
各$C_{ij}$は有限自由$B$-加群とし、同型は特殊ファイバー上で恒等とする。
有限自己同型群による粗モジュライの商と、この基点付き変形問題は区別する。

\begin{lemma}[相対Koszul性と相対4周期性]\label{lem:relative-four}
標準14641型$C_0$の成分ごとに平坦な二次変形$C_B$は相対Koszulであり、
双対の各成分は自由で特殊ファイバーと同じ階数を持ち、相補次数積は完全である。
従って$C_B$はFrobenius標準4周期構造を持つ。
trace変更の頂点スカラー共役と、同型によるtrace移送も$B$上で成立する。
\end{lemma}
\begin{proof}
内部次数$n$を固定した相対bar複体は、有界な有限自由$B$-加群の複体である。
剰余体上ではKoszul性によりホモロジーはホモロジー次数$n$にのみ存在する。
剰余体上の可縮な二項対を、単元の行列成分を使うGaussian消去で$B$上へ持ち上げて除く。
残る複体は一つの次数にのみ自由加群を持つため、ホモロジーは自由で、他の次数では零である。
これを各内部次数で行うと相対Koszul性と、対角Tor、従って二次双対の階数が得られる。
特殊ファイバーで完全な積対合は、その行列式が$B$の単元なので完全である。
最高次数成分は自由階数1だからtraceを選べる。
前節の$\nu$の構成、積の保存、逆双対化は完全対合だけを用い、$B$上でも同じである。
traceの比は$B^\times$に属し、$s_{i+4}/s_i=t_i$は四つの剰余類で解ける。
従って共役性と同型による移送も成立する。
\end{proof}

\paragraph{PBW座標。}
生成元を$x_0<x_1<x_2<x_3$とし、基点の収束二次書換え則を
\[
 x_bx_a\longmapsto F_{ba}(x_0,x_1,x_2,x_3)\quad(a<b)
\]
と書く。右辺は順序付き二次単項式10個の線型結合とする。
位置$i$での変形は$F_{ba}+\alpha_{i,ba}$で表し、$\alpha_i$には$6\cdot10=60$成分がある。
基点での書換えは単項式順序を減らし、変形部分はArtin環の極大イデアルの次数を上げる。
従ってこの二つを組み合わせた有限の還元手順により、Artin環上の正規形を定義できる。
唯一の重なりは$x_cx_bx_a$（$a<b<c$）の4個である。
各重なりの差を順序付き三次単項式20個に展開し、全$4\cdot20=80$成分を零とすることが、
相対Diamond lemmaによる平坦性の必要十分条件になる。
これは通常の数値的Hilbert一致だけをAS正則性の証明に使う議論ではない。

\paragraph{一次複体。}
基底の一次変更を$f_i\in\operatorname{End}(k^4)$とする。
関係の変化と重なりの線型項は
\begin{equation}\label{eq:linear-complex}
 \alpha_i=N_0f_i+N_1f_{i+1},\qquad
 M_0\alpha_i+M_1\alpha_{i+1}=0
\end{equation}
である。$N_0,N_1$は$60\times16$、$M_0,M_1$は$80\times60$行列である。
定義は次の通りで、同梱\path{computations/pbw_complex.py}にもそのまま実装する。
関係$r_{ba}=x_bx_a-F_{ba}$の第$j$因子を$f$で置換し基点で正規化したものが$N_jf$である。
三次重なりを左先行・右先行の二通りで還元し、位置0または1で一度だけ
$F$を$\alpha$に置換した項の差が$M_0\alpha$または$M_1\alpha$である。
全ての語・対・三つ組・単項式は辞書式順序に固定する。

周期$m\mid4$の変形をFourier分解して$\alpha_i=\zeta^i\alpha$と書くと
\begin{equation}\label{eq:fourier-complex}
 k^{16}\xrightarrow{D_1(\zeta)=N_0+\zeta N_1}
 k^{60}\xrightarrow{D_2(\zeta)=M_0+\zeta M_1}k^{80},\qquad\zeta^m=1
\end{equation}
を得る。$D_2D_1=0$は基底変更による変形が結合的であることから従う。
コードでは$M_0N_0=M_1N_1=M_0N_1+M_1N_0=0$も個別に検査する。
一次変形の同値類は$H^2_\zeta=\ker D_2(\zeta)/\operatorname{im}D_1(\zeta)$である。

\begin{lemma}[滑らかな族を用いる局所完全性]\label{lem:smooth-source}
有限型の表示空間の点$x$について、滑らかなパラメータ空間からの射が$x$で接空間を全射するなら、
$x$の表示空間は滑らかであり、その族は形式的局所的に全ての変形を含む。
同じ結論は有限変数・有限形式方程式で与えた完備局所環にも成立する。
\end{lemma}
\begin{proof}
$d=\dim T_x$とする。滑らかな源の完備局所座標から、接写像が同型となる$d$変数の形式的切片を選ぶ。
標的の完備局所環を$k[[u_1,\ldots,u_d]]/I$と書く。
切片への写像における$u_i$の像の線型部は独立なので、形式的逆写像定理により
$k[[u_1,\ldots,u_d]]\to k[[t_1,\ldots,t_d]]$は同型である。
この写像が$I$を零へ送るため$I=0$であり、結論を得る。
\end{proof}
\begin{remark}
\emph{源の滑らかさ}を省いて接写像の全射性だけを使ってはいけない。
例えば尖点曲線$y^2=x^3$の$\mathbb A^2$への包含は原点で接空間を全射するが、
曲線は平面の既約成分ではない。
本稿では\cite[Proposition 5.1]{BILN}をこの一般性では用いず、上の補題を直接使う。
\end{remark}

\section{量子多項式成分：同型類次元6と3}\label{sec:quantum}
$q_{ab}\in k^\times$（$a<b$）に対し
\[
 A_q=k\langle x_0,x_1,x_2,x_3\rangle/
 (x_bx_a-q_{ab}x_ax_b\mid a<b)
\]
と置く。$q_{ba}=q_{ab}^{-1}$と約束する。
全ての$q$で順序付き単項式が基底となり、双対量子外積代数の相補次数積が完全なので、
$A_q$は標準14641型である。次の開集合を考える。
\[
 U=\{q:c_{abc}:=q_{ab}q_{bc}/q_{ac}\ne1\quad(a<b<c)\}.
\]

\begin{theorem}[量子多項式族の二つの同型分類]\label{thm:quantum-orbits}
$U$上で次が成立する。
\begin{enumerate}
\item graded同型は、四生成元の同じ置換による$q$の変換に尽きる。
従って同型類の有理商の次元は6である。
\item 添字保持同型$\check A_q\simeq\check A_{q'}$は、上の置換に加えて
\begin{equation}\label{eq:quantum-twist}
 q'_{ab}=q_{ab}\,s_a/s_b\qquad(s_a\in k^\times)
\end{equation}
を許したものに尽きる。この同型類の次元は3である。
\item 置換で割る前の商は、例えば
$(c_{012},c_{013},c_{023})$で記述できる。
$c_{123}=c_{012}c_{023}/c_{013}$であり、独立な不変量は三つである。
\end{enumerate}
\end{theorem}
\begin{proof}
多重線型化は$p_bq_a-q_{ab}p_aq_b=0$である。
$p$の非零成分が三つ以上あると、対応する三点で$c_{abc}=1$が必要となる。
また$p$の支持と$q$の支持は一致する。従って$U$上の点スキームは六つの座標直線の和であり、
その特異点は四つの座標頂点である。関係対応は各座標頂点を同じ頂点へ送る。

この配置を保つ射影線型写像は四頂点を置換するので単項行列である。
graded同型に対しては両テンソル因子に同じ単項行列が働くため、
対角成分は$q$を変えず、置換だけが残る。

$\ZZ$-代数同型の次数一部分を$g_i$とする。
各$g_i$は単項行列であり、頂点対応の保存により隣接位置の置換は同じである。
共通の置換を除き、$g_i(x_a)=t_{i,a}x'_a$と書くと
\[
 \frac{q'_{ab}}{q_{ab}}
 =\frac{t_{i,a}t_{i+1,b}}{t_{i,b}t_{i+1,a}}.
\]
$i=0$の式だけから$q'/q$は$s_a/s_b$の形になる。
逆に式\eqref{eq:quantum-twist}があれば$t_{i,a}=s_a^{-i}$とすることで全位置に同型を作れる。
有効な対角twist群は$(k^\times)^4/k^\times$、次元3であり、自由に作用する。
その商の次元は$6-3=3$である。
三角積はこの作用で不変であり、三つの辺$q_{01},q_{02},q_{03}$を1に規格化すれば、
残りの三つの辺はこれらの三角積で復元される。
最後の$S_4$という有限群の商は次元を変えない。
\end{proof}

\paragraph{次数零自己同型とstack次元。}
一般点でgraded自己同型群は対角$(k^\times)^4$である。
従ってgraded族のstack次元と同型類次元6を混同してはならない。
$\ZZ$-代数ではさらに各頂点のスカラーがあり、ここでの3もstack次元の主張ではない。

\begin{theorem}[一般量子多項式点の局所成分]\label{thm:quantum-local}
量子多項式族のある非空開部分では、標準AS正則$\ZZ$-代数の全Artin変形は、
同型を除いて量子多項式族から得られる。
従ってこの成分の一般点で、graded同型類の局所次元は6、
裸の添字保持同型類の局所次元は3である。
この点の近傍の次元差はtwist方向の忘却によるものであり、
非1周期的対象を加えたことによる増大ではない。
この主張の行列階数部分は計算補助証明である。
\end{theorem}
\begin{proof}
次の有理数点を使う。順序は$(01,02,03,12,13,23)$である。
\begin{equation}\label{eq:quantum-sample}
 q_*=(2,3,1/6,5,2/5,15).
\end{equation}
この点では各行の積$\prod_{b\ne a}q_{ab}=1$であり、全ての三角積が1と異なる。
双対量子外積代数のNakayama作用は次数一で$-I$、最高次traceを$(-1)^i$倍すると
厳密4移動がFrobenius標準構造になる。

式\eqref{eq:fourier-complex}の厳密階数は
\[
\begin{array}{c|rrr}
\zeta&\rk D_1&\rk D_2&\dim H^2_\zeta\\\hline
1&12&42&6\\
-1&15&45&0\\
\sqrt{-1}&15&45&0\\
-\sqrt{-1}&15&45&0
\end{array}
\]
である。さらに六つの$q$の微分を$K$と書くと
$D_2(1)K=0$、$\rk(D_1(1)\mid K)=18$である。
これらは\path{computations/verify_moduli_jets.py}で$\QQ$および$\QQ(\sqrt{-1})$上で検証する。
行列の定義は前節の有限還元手順であり、入力は式\eqref{eq:quantum-sample}だけである。

四窓表示の接空間は$6+57=63$次元である。
滑らかな族$\GL_4^4\times(k^\times)^6$の接写像は、基底変更の57方向と$q$の6方向により
全63次元を張る。補題\ref{lem:smooth-source}により四周期PBW表示の形式的局所空間は滑らかで、
この族が全変形を含む。
相対Diamond lemmaにより、ここで使う形式方程式は平坦二次変形を記述する。
補題\ref{lem:relative-four}によって、裸の標準AS正則変形も標準4周期座標に直せるため、
この議論から漏れる無周期の形式的方向はない。

従って裸の同型類は量子多項式族における同型類の商で尽くされる。
前定理の三次元twist群と有限置換群による分類を適用して局所次元3を得る。
graded側には同じ議論を$\zeta=1$で適用して局所次元6を得る。
非零小行列式とKodaira--Spencer写像の階数条件は開条件なので、結論は非空開部分で成立する。
滑らかな形式的近傍は既約であるから、これは族内の次元だけでなく、対応する一般点の局所成分の次元である。
\end{proof}

\begin{remark}
この定理は全ての成分が量子多項式成分であるとの主張ではない。
特に次節のClifford点は異なる変形挙動を示す。
また一般の4周期表示の接空間の6を、裸の$\ZZ$-代数の同型類次元と読み替えてはいけない。
標準化しない周期構造と、twistによる同一視の余分な自由度が残っているためである。
\end{remark}

\section{Clifford特殊点の一次変形と二次障害}\label{sec:clifford}
式\eqref{eq:period-two-family}の$\lambda=0$を$C_0=\check A_0$とする。
ここでは$x_0=W,x_1=Y,x_2=Z,x_3=X$の順序を使う。書換えは
\begin{equation}\label{eq:clifford-rules}
\begin{array}{lll}
YW\mapsto-WY,& ZW\mapsto-WZ,& ZY\mapsto-YZ+W^2,\\
XW\mapsto-WX+Z^2,& XY\mapsto-YX+Z^2+W^2,& XZ\mapsto-ZX.
\end{array}
\end{equation}
この六規則の四つの重なりは解消する。
双対を$x,y,z,w$（それぞれ$X,Y,Z,W$に双対）で書くと
\[
 E_0=k[x,y,z,w]/(x^2,y^2,z^2+xw+xy,w^2+yz+xy)
\]
である。四つの二次式の共通零点は原点だけなので正則列であり、
$E_0$はHilbert級数$(1+t)^4$のFrobenius代数である。
元代数の二次PBW性と併せて$C_0$の標準14641型を確認できる。
$E_0$は可換なので一定traceのNakayama作用は恒等であり、厳密4移動がFrobenius標準である。

\begin{theorem}[Clifford点の変形接空間]\label{thm:clifford-tangent}
式\eqref{eq:clifford-rules}に対するFourier複体の厳密階数は
\begin{equation}\label{eq:clifford-ranks}
\begin{array}{c|rrr}
\zeta&\rk D_1&\rk D_2&\dim H^2_\zeta\\\hline
1&15&36&9\\
-1&15&39&6\\
\sqrt{-1}&15&45&0\\
-\sqrt{-1}&15&45&0
\end{array}
\end{equation}
である。選んだgraded表示$A_0$から裸の$C_0$への忘却の接写像は単射で、
その像は9次元の偶成分、余核は6次元の交代成分である。
$C_0$の基点付き標準AS正則$\ZZ$-代数変形の接空間は15次元である。
\end{theorem}
\begin{proof}
階数は式\eqref{eq:clifford-rules}から前節の算法で作った行列を有理数上で消去して得る。
同梱コードは$D_2D_1=0$、各階数、核と像の商の基底を独立に検査する。
四周期変形は四つのFourier成分の直和だから、その商接空間は$9+6=15$であり、
graded変形の像は$\zeta=1$成分である。

裸の変形への移行でこの次元を失わないことを確認する。
$C_0$の任意の無限小自己同型は、標準4周期構造の自然性により、
頂点スカラーを除けば4周期的である。
各Fourier成分で$D_1$の核はスカラー行列の一次元だけなので、
射影的な無限小自己同型は零である。
Artin変形上で恒等に還元される自己同型も、極大イデアルのべきに沿う帰納法により頂点スカラーに限る。
従って厳密4周期構造と相対標準4周期構造との差は、頂点スカラー共役で取り除ける。
二つの裸の同型の比較にも同じ議論が適用できる。
補題\ref{lem:relative-four}により全ての裸の変形を4周期表示にできるので、
ここでは4周期構造を忘れても、スカラーを除いた基点付き変形接空間は変わらない。
\end{proof}

\paragraph{座標の固定。}
各モードで$\ker D_2$と$\operatorname{im}D_1$を通常の有理数rrefで計算し、
像の列基底の後に核の列基底を並べて、ピボットとして新しく残る核の列を商基底とする。
この決定的な手順で偶基底$e_0,\ldots,e_8$、交代基底$o_0,\ldots,o_5$を固定する。
全係数は\path{research/moduli_certificate.json}に記録する。
特に$o_0$は$YW$の右辺へ$ZX$を、$o_5$は$XZ$の右辺へ$WY$を加える基底ベクトルである。
位置$i$での交代変化には$(-1)^i$を掛ける。

\begin{theorem}[全15方向の二次障害]\label{thm:clifford-quadratic}
上の基底で、一次変形$e+o$が二次まで持ち上がるための必要十分条件は
\begin{equation}\label{eq:clifford-Q}
 Q(o)=o_0o_5+o_1o_2-o_2o_3+o_3o_4=0
\end{equation}
である。偶方向同士と偶・交代の混合の二次障害は零である。
$Q$の階数は6である。従って基点付き変形のminiversal環を$R$と書くと、適当な形式座標で
\[
 R/\mathfrak m^3\simeq
 k[e_0,\ldots,e_8,o_0,\ldots,o_5]/(Q,(e,o)^3).
\]
この命題は\emph{二次まで}の主張であり、完備局所環$R$全体を決定するものではない。
\end{theorem}
\begin{proof}[計算補助証明]
書換え則を$F+t\alpha+t^2\beta$とし、三次重なりの$t^2$係数を取る。
その式は$D_2\beta+\operatorname{Ob}_2(\alpha)=0$である。
$\alpha_i=e+(-1)^i o$なので、二次残差は偶・交代成分だけを持つ。
$L_+$を$D_2(1)^{\mathsf T}$の核の転置、$L_-$を$D_2(-1)^{\mathsf T}$の核の転置とする。
それぞれ$44\times80$、$41\times80$であり、
残差が$D_2$の像に入ることは$L_+\operatorname{Ob}_{2,+}=0$と
$L_-\operatorname{Ob}_{2,-}=0$に同値である。

15個の基底ベクトルとその異なる二つの和で残差を評価する。
二次式であるため、この$15+\binom{15}{2}=120$回の厳密計算で全係数が決まる。
得られた85成分の二次式は全て$0,Q,-Q$のいずれかであり、少なくとも一つが$Q$または$-Q$である。
これを\path{computations/verify_moduli_jets.py}が有理数のみで検査する。
$\pm\sqrt{-1}$成分には残差がないので、四周期の二次補正を許しても条件は変わらない。
前定理の相対標準化により裸の変形でも同じ条件である。

この検査は単に正の持上げ候補を見つけるものではない。
$L_\pm$が左核全体の基底なので、像への所属の必要性と十分性を同時に証明している。
また一般の二次式の全係数を検査しており、特定のパラメータ値だけの数値観察ではない。
$rref$によって線型方程式と基底変更方向を除いた局所方程式の二次部分も同じ残差であるため、
表示した$R/\mathfrak m^3$を得る。
最後に
\[
 z_1=o_0,\ z_2=o_5,\ z_3=o_1-o_3,\ z_4=o_2,\ z_5=o_3,\ z_6=o_4
\]
は可逆な線型変換であり、$Q=z_1z_2+z_3z_4+z_5z_6$となる。従って階数は6である。
\end{proof}

\begin{corollary}[Pl\"ucker型の最初の障害]\label{cor:plucker-obstruction}
六つの交代方向に対する二次障害の零点錐は、
Pl\"ucker埋込み$\Gr(2,4)\subset\mathbb P^5$上のアフィン錐と線型同型である。
これは\emph{二次障害が許す候補錐}の記述であり、実際の接錐や成分全体の同定ではない。
\end{corollary}
\begin{proof}
$\Gr(2,4)$のPl\"ucker式$p_{01}p_{23}-p_{02}p_{13}+p_{03}p_{12}=0$に
$(p_{01},p_{23},p_{02},p_{13},p_{03},p_{12})=(z_1,z_2,z_3,-z_4,z_5,z_6)$を代入する。
\end{proof}

\begin{proposition}[二つの積分可能な軸と混合障害]\label{prop:two-axes}
式\eqref{eq:clifford-rules}のうち、それぞれ一方だけを
\[
 YW\mapsto-WY+(-1)^i aZX,\qquad
 XZ\mapsto-ZX+(-1)^i bWY
\]
に置き換えると、各軸上で全パラメータの二次PBW族を得る。
原点での形式変形は相対AS正則である。
しかし両軸の接方向を同時に取った一次変形には二次障害$ab$があり、
$a=b=1$という接方向は、任意の二次補正を許しても持ち上がらない。
\end{proposition}
\begin{proof}
各軸について二つの始点と四つの三次重なり、合計8個の重なり差は$\QQ[a]$または$\QQ[b]$上で恒等的に零である。
同梱コードは記号多項式としてこれを確認する。
$b$軸では通常の次数付き辞書式順序を使える。
$a$軸では重み$(W,Y,Z,X)=(1,4,1,1)$を先に比較し、同じ重みでは$X>Z>Y>W$の辞書式順序を使うと
全規則が減少する。Diamond lemmaから順序付き単項式基底と平坦性を得る。
形式的なAS正則性は補題\ref{lem:relative-four}による。
新しい$a$軸の\emph{全閉ファイバー}のAS正則性は、この議論だけでは主張しない。
接方向は$o_0,o_5$なので、式\eqref{eq:clifford-Q}を制限すると$ab$になる。
\end{proof}

\subsection{まだ結論できないこと}
上の計算から、Clifford点における成分次元が14であるとは結論できない。
三次以上の独立な関係式が残る可能性がある。
また$9\to15$は\emph{選んだgraded表示}$A_0$での忘却写像の微分である。
同じ裸の$C_0$に別の1周期構造を付けたときのgraded表示は別途調べる必要がある。
全ての1周期的部分の局所余次元を6または5と断定してはならない。

\begin{conjecture}[検証対象の局所模型]\label{conj:clifford-model}
Clifford点の基点付き変形環は
\[
 k[[e_0,\ldots,e_8,z_1,\ldots,z_6]]/
 (z_1z_2+z_3z_4+z_5z_6)
\]
と同型である。
\end{conjecture}
この予想を裏付ける確定事項は定理\ref{thm:clifford-quadratic}の二次切断までである。
この模型が正しければ局所次元14となるが、その十分性・高次持上げ・全周期1構造の分類は未解決である。

